\chapter{Football Model with SQMC}
\label{ch:model}
In this section, we demonstrate how limitations of SQMC can be overcome by performing Rao-Blackwellization and using SQMC on a smaller effective dimension of the parameter space.

% ADD IN CURSE OF DIMENSIONALITY FROM \citp{chopingerber2017}

\section{Background}
In a general state-space model, the joint distribution of latent states and observations can be factorized as 

\begin{equation}
p(x_{0:T}, y_{1:T}) = p(x_0) \prod_{t=1}^T p(x_t|x_{t-1})p(y_t|x_t) \tag{2.1}
\end{equation}

where $x_t$ is the latent state at time $t$ and $y_t$ is the observation at time $t$. The assumptions in this model are that latent states $x_t$ are Markovian and that observations $y_t$ are conditionally independent given the latent states. 

Duffield, Power and Rimella \citep{duffieldpowerrimella2024factorialssm} introduces the concept of a factorial state-space Model (fSSM), which extends a general state-space model by making the following assumptions: (1) the initial latent states of index $k$ are independent (2) evolution of latent states of index $k$, $p(x_t \mid x_{t-1})$ are independent and (3) given the latent states $x_t$, the observations $y_t$ are conditionally independent of all other indexes $-k$. For independent latent factors with a likelihood that depends on a local subset $\mathcal{O}_t$ of the latent indexes involved in observation $t$, the joint distribution of the fSSM is given by

\begin{equation}
p(x_{0:T}^{1:K}, y_{1:T}) = \left[ \prod_{k=1}^K p(x_0^k) \prod_{t=1}^T p(x_t^k|x_{t-1}^k) \right] \cdot \prod_{t=1}^T p(y_t|x_t^{\mathcal{O}_t}) \tag{2.2}
\end{equation}

where $x_t^k$ is the latent state of index $k$ at time $t$ and $\mathcal{O}_t$ is the set of indexes involved in observation $t$; for a football match $\mathcal{O}_t = \{i_t, j_t\}$. This factorization is extremely useful as it reduces the computation of a state-space model from dimension $\mathcal{O}(N \times K)$ terms to $\mathcal{O}(N + K)$ terms, where $N$ is the number of particles and $K$ is the number of indexes.

The key assumption in this model is that latent states of index $k$ are independent. However, in many applications, this assumption is not valid. In the context of football, the evolution of national teams are not independent over time because of the influence of players from different nationalities playing in the same league or club. We therefore emphasise that the football model we develop below is not a factorial state-space model under this independence definition when $\Gamma_0$ is non-diagonal: it is a correlated Gaussian state-space model with a sparse likelihood, or equivalently a relaxation of the factorial model. To overcome the limitation of independence, we demonstrate that we can perform Rao-Blackwellization on the latent states of index $k$. We also demonstrate how we can use SQMC on a smaller effective dimension of parameter space to improve the performance of the model.

\section{Model}
\label{sec:model}
We attempt to model the performance of teams during the World Cup 2026. There are $K = 48$ teams competing in the World Cup 2026. We represent the latent states of each team as a 2-dimensional vector $x = (x^{att}, x^{def}) \in \mathbb{R}^2$, which represents the attacking and defensive strengths of a national team. The full latent state is therefore $x \in \mathbb{R}^{2K}$, and our initial distribution of latent states is given by 

\begin{equation}
    x_0 \in \mathbb{R}^{2K}, \qquad x_0 \sim \mathcal{N}_{2K}(\mu_0, \Sigma_0)
\end{equation}

where $\mu_0 = \mathbf{0}_{2K}$ for simplification. We make the further assumption that the relationship between attack and defense of all teams are the same, represented by the matrix $B \in \mathbb{R}^{2 \times 2}$. This assumption allows us to simplify the covariance matrix into a Kronecker product of the form $\Sigma_0 = \Gamma_0 \otimes B$, where $\Gamma_0 \in \mathbb{R}^{K \times K}$ is a covariance matrix that represents the relationship between teams. We require the conditions $\Gamma_0 \succ 0$, $B \succ 0$ and $\det(B) = 1$; the last condition removes the scale non-identifiability of the Kronecker parameterization, since $(c\Gamma_0) \otimes (B/c) = \Gamma_0 \otimes B$. We model the evolution of latent states as an Ornstein-Uhlenbeck process, a mean-reverting stochastic process, given by

\begin{equation}
    x_t \mid x_{t-1} \sim \mathcal{N}(\mu_0 + \phi_t (x_{t-1} - \mu_0), Q_t)
\end{equation}

where we define $\Delta t_t = t_t - t_{t-1}$, $\phi_t = \exp(-\kappa \Delta t_t)$ and $Q_t = (1 - \phi_t^2)\Sigma_0$ is the covariance matrix of the process noise. $\kappa$ is the mean-reversion rate and $\Delta t_t$ is the time difference between $t$ and $t-1$, measured in days. This behaviour is desirable in football as we expect average skill levels to revert to a long-run marginal distribution of $\mathcal{N}(\mu_0, \Sigma_0)$ over time without any new observations.

Since each match only involves two teams, we represent the latent states of the two playing teams as $x_t^{\mathcal{O}_t} = (x_t^{(i_t)}, x_t^{(j_t)})$. Our observation is a bivariate vector of scores for team $i_t$ and team $j_t$, represented as $y_t = (y_t^{(i_t)}, y_t^{(j_t)})$. The likelihood function follows a bivariate poisson distribution introduced in \citep{karlisntzoufras2003bivariatepoisson} for sports data.

\begin{equation}
    G_t(y_t \mid x_t^{\mathcal{O}_t}) = e^{-(\lambda_1 + \lambda_2 + \lambda_3)} \frac{\lambda_1^{y_t^{(i_t)}}}{y_t^{(i_t)}!} \frac{\lambda_2^{y_t^{(j_t)}}}{y_t^{(j_t)}!} \sum_{k=0}^{\min(y_t^{(i_t)}, y_t^{(j_t)})} \binom{y_t^{(i_t)}}{k} \binom{y_t^{(j_t)}}{k} k! \left( \frac{\lambda_{3}}{\lambda_1 \lambda_2} \right)^k
\end{equation}

where $\lambda_1 = \exp(\alpha + x_t^{\text{att}, i_t} - x_t^{\text{def}, j_t})$, $\lambda_2 = \exp(\alpha + x_t^{\text{att}, j_t} - x_t^{\text{def}, i_t})$, $\lambda_3 = \exp(\beta)$. The parameter $\alpha$ is a common log scoring-rate intercept because it enters both $\lambda_1$ and $\lambda_2$ symmetrically; it is not a home-advantage parameter. The parameter $\beta$ is the log intensity of the shared Poisson component, $\lambda_3 = \exp(\beta)$; it controls dependence but is not itself the score correlation.

\section{Rao-Blackwellized SMC (RB-SMC)}
\label{sec:rb-smc}
In section~\ref{sec:smc}, we describe the SMC algorithm for a general state-space model. In our previous section~\ref{sec:model}, we state that only the latent states of the two playing teams $x_t^{\mathcal{O}_t}$ are relevant to the likelihood function $G_t$. The RB-SMC targets the marginal posterior of the sequence of sampled coordinates and recovers the non-playing teams in closed form instead. Defining the time-varying non-playing latent states as $x_t^{\mathcal{R}_t} = x_t^{-\mathcal{O}_t}$, and interpreting $x_{0:T}^{\mathcal{O}}$ in the path factorization as the time-varying collection $\{x_t^{\mathcal{O}_t}\}_{t=0}^T$, we can factorize the joint posterior as follows,

\begin{equation}
    p(x_{0:T} \mid y_{1:T}) = p(x_{0:T}^{\mathcal{O}} \mid y_{1:T}) \cdot p(x_{0:T}^{\mathcal{R}} \mid x_{0:T}^{\mathcal{O}}, y_{1:T})
\end{equation}

Since the likelihood depends only on $x_{0:T}^{\mathcal{O}}$, the second factor simplifies to
\begin{equation}
    p(x_{0:T}^{\mathcal{R}} \mid x_{0:T}^{\mathcal{O}}, y_{1:T}) = p(x_{0:T}^{\mathcal{R}} \mid x_{0:T}^{\mathcal{O}}).
\end{equation}

Since our transition model is linear Gaussian, the covariance recursions are deterministic and do not depend on the sampled coordinates or the particle paths; hence $\Sigma_t$ is shared across all particles while the means $\mu_t^{(i)}$ carry the particle-specific information. Under the bootstrap proposal, the RB-SMC algorithm is given by the following steps,

\textbf{1 Prediction}: Compute the predictive mean and covariance matrix for the full state,

\begin{equation}
    \mu_{t|t-1}^{(i)} = \mu_0 + \phi_t (\mu_{t-1}^{(i)} - \mu_0) \nonumber
\end{equation}

\begin{equation}
    \Sigma_{t|t-1} = \phi_t^2 \Sigma_{t-1} + (1 - \phi_t^2) \Sigma_0 \nonumber
\end{equation}

\textbf{2 Bootstrap Proposal Sampling}: For particle $i$, the full predictive distribution is
\begin{equation}
x_t \mid I_{t-1}=i,y_{1:t-1} \sim \mathcal{N}\left(\mu_{t\mid t-1}^{(i)}, \Sigma_{t\mid t-1} \right). \end{equation}

At each observation, only the two playing teams are sampled. We extract the predictive mean and covariance for the playing teams' coordinates $x_t^{\mathcal{O}_t}$, and the following predictive marginal under the bootstrap proposal is defined as,

\begin{equation}
    q_t^{(i)}(x_t^{\mathcal{O}_t}) = p(x_t^{\mathcal{O}_t} \mid I_{t-1} = i, y_{1:t-1}) = \mathcal{N}(\mu_{t|t-1}^{\mathcal{O}_t,(i)}, \Sigma_{t|t-1}^{\mathcal{O}_t\mathcal{O}_t}) \nonumber
\end{equation}

\textbf{3 Computing log-weights}: For the propagate-weight-resample ordering, compute the unnormalized log-weights for each particle $(i)$ using the previous weight and the likelihood function $G_t$,

\begin{equation}
    \log \tilde{w}_t^{(i)} = \log w_{t-1}^{(i)} +  \log G_t(y_t \mid x_t^{\mathcal{O}_t, (i)}) \nonumber
\end{equation}

and normalize,
\begin{equation}
    w_t^{(i)} = \frac{\tilde{w}_t^{(i)}}{\sum_{n=1}^N \tilde{w}_t^{(n)}}. \nonumber
\end{equation}

\textbf{4 Exact Marginalization of Non-Playing Teams}: Given the predictive joint distribution and conditioning on the playing teams' coordinates $x_t^{\mathcal{O}_t}$, the non-playing coordinates $x_t^{\mathcal{R}_t}$ are Gaussian with a closed-form mean and covariance matrix defined by the Kalman filter equations,

\begin{equation}
    x^{\mathcal{R}_t,(i)}_{t \mid t-1} \mid x^{\mathcal{O}_t,(i)}_{t} \sim \mathcal{N}\left( \mu^{\mathcal{R}|\mathcal{O},(i)}_{t \mid t-1}, \Sigma^{\mathcal{R}|\mathcal{O}}_{t \mid t-1} \right) \nonumber
\end{equation}

\begin{align}
    \mu^{\mathcal{R}_t|\mathcal{O}_t,(i)}_{t \mid t-1} &= \mu_{t|t-1}^{\mathcal{R}_t} + K_t \left( x_t^{\mathcal{O}_t,(i)} - \mu_{t|t-1}^{\mathcal{O}_t} \right), \nonumber \\
    \Sigma^{\mathcal{R}_t|\mathcal{O}_t}_{t \mid t-1} &= \Sigma_{t|t-1}^{\mathcal{R}_t\mathcal{R}_t} - K_t \Sigma_{t|t-1}^{\mathcal{O}_t\mathcal{R}_t}, \nonumber \\
    K_t &= \Sigma_{t|t-1}^{\mathcal{R}_t\mathcal{O}_t} \left( \Sigma_{t|t-1}^{\mathcal{O}_t\mathcal{O}_t} \right)^{-1}. \nonumber
\end{align}

In section~\ref{sec:smc}, we made the assumption that our initial covariance is a Kronecker product of the form $\Sigma_0 = \Gamma_0 \otimes B$. Since our transition model is linear Gaussian, the Kronecker structure is preserved recursively by induction, 

\begin{align}
    \Sigma_{t \mid t-1} &= \phi_t^2 \Sigma_{t-1} + (1 - \phi_t^2) \Sigma_0 \nonumber \\
    &= \phi_t^2 (\Gamma_{t-1} \otimes B) + (1 - \phi_t^2) (\Gamma_0 \otimes B) \nonumber \\
    &= \left( \phi_t^2 \Gamma_{t-1} + (1 - \phi_t^2) \Gamma_0 \right) \otimes B \nonumber \\
    &= \Gamma_{t|t-1} \otimes B. \nonumber
\end{align}

so that the covariance blocks for the playing and non-playing teams factor as

\begin{equation}
    \Sigma_{t|t-1}^{\mathcal{O}_t\mathcal{O}_t} = \Gamma^{\mathcal{O}_t\mathcal{O}_t}_{t|t-1} \otimes B, \qquad
    \Sigma_{t|t-1}^{\mathcal{R}_t\mathcal{R}_t} = \Gamma^{\mathcal{R}_t\mathcal{R}_t}_{t|t-1} \otimes B, \qquad
    \Sigma_{t|t-1}^{\mathcal{O}_t\mathcal{R}_t} = \Gamma^{\mathcal{O}_t\mathcal{R}_t}_{t|t-1} \otimes B \nonumber
\end{equation}

where $\Gamma^{\mathcal{O}_t\mathcal{O}_t}_{t|t-1} \in \mathbb{R}^{2 \times 2}$ is the submatrix of $\Gamma_{t|t-1}$ for the two playing teams and $\Gamma^{\mathcal{R}_t\mathcal{R}_t}_{t|t-1} \in \mathbb{R}^{(K-2) \times (K-2)}$ is the submatrix for the non-playing teams. Using the mixed-product property $(A \otimes B)(C \otimes D) = (AC) \otimes (BD)$ and the inverse property $(A \otimes B)^{-1} = A^{-1} \otimes B^{-1}$, the Kalman gain factorizes as

\begin{align}
    K_t &= \Sigma_{t|t-1}^{\mathcal{R}_t\mathcal{O}_t} \left( \Sigma_{t|t-1}^{\mathcal{O}_t\mathcal{O}_t} \right)^{-1} \nonumber \\
    &= \left( \Gamma^{\mathcal{R}_t\mathcal{O}_t}_{t|t-1} \otimes B \right) \left( \Gamma^{\mathcal{O}_t\mathcal{O}_t}_{t|t-1} \otimes B \right)^{-1} \nonumber \\
    &= \left( \Gamma^{\mathcal{R}_t\mathcal{O}_t}_{t|t-1} \otimes B \right) \left( (\Gamma^{\mathcal{O}_t\mathcal{O}_t}_{t|t-1})^{-1} \otimes B^{-1} \right) \nonumber \\
    &= \left( \Gamma^{\mathcal{R}_t\mathcal{O}_t}_{t|t-1} (\Gamma^{\mathcal{O}_t\mathcal{O}_t}_{t|t-1})^{-1} \right) \otimes \left( B B^{-1} \right) \nonumber \\
    &= \tilde{K}_t \otimes I_2 \nonumber
\end{align}

\begin{equation}
    \tilde{K}_t = \Gamma^{\mathcal{R}_t\mathcal{O}_t}_{t|t-1} (\Gamma^{\mathcal{O}_t\mathcal{O}_t}_{t|t-1})^{-1} \nonumber
\end{equation}

and the conditional covariance of the non-playing teams becomes

\begin{equation}
    \Sigma_{t|t-1}^{\mathcal{R}_t|\mathcal{O}_t} = \left( \Gamma^{\mathcal{R}_t\mathcal{R}_t}_{t|t-1} - \tilde{K}_t \Gamma^{\mathcal{O}_t\mathcal{R}_t}_{t|t-1} \right) \otimes B \nonumber
\end{equation}

To preserve the Kronecker structure after conditioning, we update the full component mean and the full covariance. Define the team-space Kalman gain

\begin{equation}
    \overline{K}_t = \Gamma_{t|t-1}^{:\mathcal{O}_t} \left( \Gamma_{t|t-1}^{\mathcal{O}_t\mathcal{O}_t} \right)^{-1}, \nonumber
\end{equation}

where $\Gamma_{t|t-1}^{:\mathcal{O}_t}$ contains all rows and the two playing-team columns of $\Gamma_{t|t-1}$. The complete component mean is updated as

\begin{equation}
    \mu_t^{(i)} = \mu_{t|t-1}^{(i)} + (\overline{K}_t \otimes I_2) \left( x_t^{\mathcal{O}_t,(i)} - \mu_{t|t-1}^{\mathcal{O}_t,(i)} \right). \nonumber
\end{equation}

The rows of $\overline{K}_t$ belonging to $\mathcal{O}_t$ form the identity, so the playing-team entries of $\mu_t^{(i)}$ equal the sampled entries $x_t^{\mathcal{O}_t,(i)}$, while the $\mathcal{R}_t$ rows recover the conditional mean of the non-playing teams in closed form. The full filtered covariance is defined by the Schur complement

\begin{equation}
    \Gamma_t = \Gamma_{t|t-1} - \overline{K}_t \Gamma_{t|t-1}^{\mathcal{O}_t} = \Gamma_{t|t-1} - \Gamma_{t|t-1}^{:\mathcal{O}_t} \left( \Gamma_{t|t-1}^{\mathcal{O}_t\mathcal{O}_t} \right)^{-1} \Gamma_{t|t-1}^{\mathcal{O}_t} \nonumber
\end{equation}

and hence
\begin{equation}
    \Sigma_t = \Gamma_t \otimes B. \nonumber
\end{equation}

This full update encodes both effects of conditioning in a single matrix: the rows and columns of $\Gamma_t$ corresponding to the playing teams $\mathcal{O}_t$ are zero, because their coordinates are now known (sampled) values with no remaining uncertainty, while the $\mathcal{R}_t\mathcal{R}_t$ block of $\Gamma_t$ is the nonzero conditional covariance of the non-playing teams, which is precisely the Schur complement $\Sigma_{t|t-1}^{\mathcal{R}_t|\mathcal{O}_t}$ derived above.

The Kronecker structure provides computational benefits for the algorithm. Instead of computing Kalman recursions on the full $2K \times 2K$ covariance matrix $\Sigma$, our Kalman recursions are computed on the $K \times K$ between-team covariance matrix $\Gamma$ and the $2 \times 2$ matrix $B$. The covariance recursion is shared across all particles, so we only need to compute it for $T$ time steps, rather than $N \times T$ times for all particles. With dense storage, updating the covariance costs $\mathcal{O}(K^2)$ per observation after a $2 \times 2$ solve, while updating all particle means costs $\mathcal{O}(NK)$. Thus the method reduces constants and avoids particle-specific covariance recursions.

\textbf{5 Resampling}: Resample the particles according to the normalized weights $\{w_t^{(i)}\}_{i=1}^N$ and set the weights to $w_t^{(i)} = 1/N$ for all particles.

\section{RB-SMC Smoothing}
\label{sec:rb-smc-smoothing}

In standard particle-based Forward Filtering Backward Simulation (FFBSi), the backward state at time $t$ is sampled from the discrete set of filtering particles, with probabilities determined by their filtering weights and compatibility with the sampled state at time $t+1$ \citep{godsilldoucetwest2004montecarlosmoothing}. 

% In our RB-SMC formulation we cannot perform backward simulation by drawing samples one step at a time from the filter results (as in the standard FFBSi algorithm), because the complementary coordinates $x_t^{\mathcal{R}}$ are marginalized analytically and their conditional distribution depends on the entire observed-coordinate trajectory $x_{1:t}^{\mathcal{O}}$, not just $x_t^{\mathcal{O}}$. Thus we need to sample the whole trajectories from the sample histories produced by the RB-SMC algorithm, similar to the case in \citep{sarkkabunchgodsill2012rbpsmoother}.

In our RB-SMC formulation, discrete point-particle FFBSi cannot be applied directly, because each filtering particle represents a Gaussian component rather than a complete state. Nevertheless, a one-step backward kernel is available in closed form: first select a Gaussian component using its backward weight, and then draw the full state from that component's conditional Gaussian distribution.

Since our transition is linear and Gaussian, our filtering distribution can be approximated as a weighted Gaussian mixture with a shared covariance matrix $\Sigma_t$ across all particles. 

\begin{equation}
    p(x_{T} \mid y_{1:T}) \approx \sum_{i=1}^N w_{T}^{(i)} \mathcal{N}(\mu_T^{(i)}, \Sigma_T)
\end{equation}

The algorithm then follows these steps,

\textbf{1 Initialization}: Sample the terminal particle index $I_T$ using the normalized weights $w_T^{(i)}$, and draw the full terminal state from the corresponding Gaussian mixture,
\begin{equation}
    I_T \sim \operatorname{Categorical}(w_T^{1:N}), \qquad X_T^* \sim \mathcal{N}(\mu_T^{(I_T)}, \Sigma_T). \nonumber
\end{equation}

\textbf{2 Backward Sampling}: For $t = T-1, \dots, 0$, sample the particle index $I_t$ from the backward kernel $p(X_t \mid X_{t+1}^*, y_{1:t})$ and draw the full state from the corresponding Gaussian mixture. By Gaussian conjugacy, the backward kernel is itself a Gaussian mixture,
\begin{equation}
    p(X_t \mid X_{t+1}^*, y_{1:t}) \approx \sum_{i=1}^N w_{t \mid t+1}^{(i)} \mathcal{N}(\mu_{t \mid t+1}^{(i)}, \Sigma_{t \mid t+1}), \nonumber
\end{equation}
where the backward weights are
\begin{equation}
    w_{t \mid t+1}^{(i)} \propto w_t^{(i)} \mathcal{N}\left( X_{t+1}^* \mid a_{t+1}^{(i)}, R_{t+1} \right), \nonumber
\end{equation}
using the predictive mean and covariance
\begin{equation}
    a_{t+1}^{(i)} = \mu_0 + \phi_{t+1}(\mu_t^{(i)} - \mu_0), \qquad R_{t+1} = \phi_{t+1}^2 \Sigma_t + (1 - \phi_{t+1}^2) \Sigma_0, \nonumber
\end{equation}
and, by the Rauch-Tung-Striebel (RTS) equations with the gain
\begin{equation}
    J_t = \Sigma_t \phi_{t+1} R_{t+1}^{-1}, \nonumber
\end{equation}
\begin{align}
    \mu_{t \mid t+1}^{(i)} &= \mu_t^{(i)} + J_t \left( X_{t+1}^* - a_{t+1}^{(i)} \right), \nonumber \\
    \Sigma_{t \mid t+1} &= \Sigma_t - J_t R_{t+1} J_t^T. \nonumber
\end{align}

One complete smoothed trajectory is obtained by one backward pass; the pass must be repeated independently to obtain $M$ trajectories. The resulting algorithm is similar to the pure particle smoother of \citet{godsilldoucetwest2004montecarlosmoothing}, but the backward weights combine the forward filtering weights with the predictive density, while the RTS backward estimates determine the Gaussian mixture component from which the full state is drawn \citep{sarkkabunchgodsill2012rbpsmoother}. This process is extremely parallelizable due to the independence of each smoothed trajectory.

\section{Rao-Blackwellized SQMC (RB-SQMC)}
\label{sec:rb-sqmc}

In our model, the likelihood at each observation only depends on the attack and defense strengths of the two playing teams. Rao-Blackwellization analytically marginalizes the particles of the non-playing teams, so only 4 coordinates are sampled.

We Hilbert sort the full $2K$-dimensional particle using the 4-dimensional predicted means of the two playing teams, $\Lambda_t(\mu_{t \mid t-1}^{(i)})$. Once an ancestor is selected, its complete components, is retained and propagated via the OU-process. This is analagous to step 1 in section~\ref{sec:rb-smc}. The potential function condition of \citet{chopingerber2017sqmcintroduction} is satisfied because $G_t(x_{t-1}, x_t) = G_t(x_t^{\mathcal{O}_t})$ so the likelihood does not depend on the previous state outside the propagated coorediantes. However, the complete Rao-Blackwellized transition is not determined by $\Lambda(\mu_{t-1}^{(i)})$ alone, because particle-specific non-playing means remain relevant to the succssive component. The 4-dimensional Hilbert ordering is therefore a projected dimension-reduction heuristic rather than an exact application of their transition-sufficiency result.

% In our model, the likelihood at each observation depends only on the attacking and defensive strengths of the two teams involved in the match. Rao-Blackwellization allows the remaining coordinates to be marginalized analytically, so only four latent coordinates are required to be propagated via sampling. This enable us to leverage the benefits of SQMC in the SMC operation.

% We Hilbert-sort the Rao-Blackwellized components using the four-dimensional state of the teams involved in the current match. The complete ancestor component, including the non-playing team means, is nevertheless retained and propagated. This projected ordering does not satisfy the exact transition-sufficiency condition of Chopin and Gerber (2017), so we treat it as a dimension-reduction heuristic.

% Our likelihood function satisfies the condition where $G_t$ depends only on $X_t$ and $\Lambda(x_{t - 1})$, where $\Lambda(x_{t - 1}) = x^{\mathcal{O},(i)}_{t - 1}$.

% We approximate the conditional distribution of non-playing teams as depending on the current playing teams state. This does not satisfy the exact transition-sufficiency condition used by Chopin and Gerber (2017), because non-playing component means remain relevant to later transitions. It is therefore treated as a dimension-reduction heuristic.

% The effective propagation dimension is therefore $d= 2 \times 2 = 4$, rather than the full state dimension $2K$; the full covariance matrix has dimension $2K \times 2K$. Including the additional coordinate used for ancestor selection, each SQMC point has dimension $d+1=5$. Our algorithm, the RB-SQMC, follows these steps,

% This dimension-reduction strategy follows the application described in Section~3.1 of \citet{chopingerber2017sqmcintroduction} and helps preserve the low-discrepancy properties of the QMC point set. 

\textbf{1 Generate N QMC points}: Generate $N$ QMC or RQMC points $(u_t^{1:N}, v_t^{1:N})$ of dimension $1 + (2 \times 2) = 5$, where $u_t$ is the first component, used for ancestor selection, and $v_t$ is the vector of $d = 4$ components used to propagate the latent states of the two playing teams, $x_t^i$ and $x_t^j$.

\textbf{2 (Projected) Hilbert Sort}: Apply the deterministic OU prediction to each Rao-Blackwellized component and extract the predicted means of the two playing teams. Define the 4-dimensional sorting key as $z_t^{(i)} = \operatorname{vec}\left( \mu_{t\mid t-1}^{\mathcal{O}_t,(i)} \right) \in \mathbb{R}^4$. Let $\psi_t:\mathbb{R}^4\rightarrow[0,1]^4$ standardize and map these coordinates to the unit hypercube, and let $h:[0,1]^4\rightarrow[0,1]$ denote the inverse Hilbert-curve mapping. Find a permutation $\sigma_t$ such that
\begin{equation}
    h\left(\psi_t(z_t^{(\sigma_t(1))})\right) \leq \cdots \leq h\left(\psi_t(z_t^{(\sigma_t(N))})\right). \nonumber
\end{equation}
Use the previous normalized particle weights in the same order,
\begin{equation}
    \left( W_{t-1}^{(\sigma_t(1))}, \ldots, W_{t-1}^{(\sigma_t(N))} \right). \nonumber
\end{equation}

Once an ancestor is selected, its complete Rao-Blackwellized component, including the means of the non-playing teams, is retained and propagated.

% For the standard SQMC justification, sort the full Rao-Blackwellized particle state, which is the full component mean $\mu_{t-1}^{(i)} \in \mathbb{R}^{2K}$. Let $\psi: \mathbb{R}^{2K} \to [0,1]^{2K}$ standardize and quantize the state, and let $h(\cdot)$ denote the Hilbert-curve index. Find the permutation $\sigma_t$ that sorts the previous $N$ particles by their Hilbert index,
% \begin{equation}
%     h\big(\psi(\mu_{t-1}^{(\sigma_t(1))})\big) \le h\big(\psi(\mu_{t-1}^{(\sigma_t(2))})\big) \le \cdots \le h\big(\psi(\mu_{t-1}^{(\sigma_t(N))})\big), \nonumber
% \end{equation}
% and use the previous normalized weights in sorted order $W_{t-1}^{(\sigma_t(1:N))}$. We use a four-dimensional projected Hilbert ordering based on the teams involved in the current match. 

% This does not satisfy the exact transition-sufficiency condition used by Chopin and Gerber (2017), because non-playing component means remain relevant to later transitions. It is therefore treated as a dimension-reduction heuristic.

\textbf{3 Resampling by inverse transform}: Sort the complete QMC pairs $(u_t^{(i)}, v_t^{(i)})$ by their first coordinate so that $u_t^{(1)} \le \cdots \le u_t^{(N)}$, preserving the pairing between each $u_t^{(i)}$ and $v_t^{(i)}$. Perform ancestor selection using the Hilbert-ordered weights $W_{t-1}$ and the first coordinate $u_t$,
\begin{equation}
    a_i = \min \left\{ j : W^{(j)} \ge u_t^{(i)} \right\}, \qquad W^{(j)} = \sum_{k=1}^{j} W_{t-1}^{(\sigma_t(k))}, \qquad b_i = \sigma_t(a_i). \nonumber
\end{equation}

\textbf{4 Propagate}: Propagate each ancestor $b_i$ through the deterministic transform of the OU-process
\begin{equation}
    x_t^{\mathcal{O}_t,(i)} = \mu_{t|t-1}^{\mathcal{O}_t,(b_i)} + L_t\, \Phi^{-1}\left( v_t^{(i)} \right), \nonumber
\end{equation}
where $L_t$ is the Cholesky factor of $\Sigma_{t|t-1}^{\mathcal{O}_t\mathcal{O}_t} = \Gamma_{t|t-1}^{\mathcal{O}_t\mathcal{O}_t} \otimes B$ and $\Phi^{-1}$ is the coordinate-wise standard normal quantile. Recover the non-playing teams' coordinates in closed form via the OU prediction and the Kalman conditioning of Section~\ref{sec:model}.

\textbf{5 Compute log-weights}: Because the ancestor weights were already used in the inverse-transform resampling, compute the unnormalized log-weights for each particle $(i)$ using only the likelihood function $G_t$,
\begin{equation}
    \log \tilde{w}_t^{(i)} = \log G_t\left( y_t \mid x_t^{\mathcal{O}_t,(i)} \right), \nonumber
\end{equation}
and normalize,
\begin{equation}
    w_t^{(i)} = \frac{\tilde{w}_t^{(i)}}{\sum_{n=1}^N \tilde{w}_t^{(n)}}. \nonumber
\end{equation}

RB-SQMC Smoothing can be performed with the same algorithm as described in Section~\ref{sec:rb-smc-smoothing}, but with the random draws replaced by QMC points. At each backward step $t$, generate $M$ QMC points $(u_t^{1:M}, v_t^{1:M})$ of dimension $1 + 2K$, sort them by the conditioning state $X_{t+1}^{*(m)}$ to preserve low-discrepancy, and propagate analytically by selecting the Gaussian component via inverse transform of the backward weights and drawing
\begin{equation}
    X_t^{*(m)} = \mu_{t \mid t+1}^{(i)} + L_{t \mid t+1}\, \Phi^{-1}\left( v_t^{(m)} \right), \nonumber
\end{equation}
where $L_{t \mid t+1}$ is the Cholesky factor of the backward covariance $\Sigma_{t \mid t+1}$.

% \section{RB-SQMC Smoothing}
% \label{sec:rb-sqmc-smoothing}

\section{Parameter Estimation}
\label{sec:parameter-estimation}

Our static parameters for the model are $\Theta = (\Gamma_0, B, \kappa, \alpha, \beta)$ where $\Gamma_0 \in \mathbb{R}^{K \times K}$, $B \in \mathbb{R}^{2 \times 2}$, $\kappa > 0$, $\alpha \in \mathbb{R}$ and $\beta \in \mathbb{R}$, subject to $\Gamma_0 \succ 0$, $B \succ 0$ and $\det(B) = 1$.

The smoothing algorithm requires the backward-sampled trajectories to be sufficiently diverse to provide a good estimate of the smoothing distribution. However, drawing many independent backward trajectories is computationally expensive: each trajectory is sampled sequentially from time $T$ to $0$, and at every time step $t$ the backward weights must be evaluated against all $N$ particles. Producing $M$ trajectories over $T$ time steps therefore costs $\mathcal{O}(MNT)$ particle-weight evaluations. Moreover, if the forward filter degenerates, the backward weights concentrate on only a few particles, so the effective number of distinct trajectories is small. The resulting smoothed trajectories then carry little diversity, which biases the estimate of the initial between-team covariance $\Gamma_0$: the backward paths are drawn from too few ancestors and the inferred spread is unreliable.

A simpler and more efficient approach is to estimate the gradient of the marginal log-likelihood. By Fisher's identity \citep{delmoral2004feynman,delmoraldoucetsingh2010forwardsmoothing}, the marginal log-likelihood of the particle filter takes the form

% We assume that the standard conditions for exiting convergence results (Moral, 2004) are satisfied, specifically specifically for consistency of estimators of the log-marginal likelihood and
% expectations under the posterior, and for unbiasedness
% of estimators of expectations under the unnormalized
% posterior.

\begin{equation}\nabla_\Theta \log p(y_{1:T} \mid \Theta) = \int p(x_{0:T} \mid y_{1:T}, \Theta) \nabla_\Theta \log p(x_{0:T}, y_{1:T} \mid \Theta) dx_{0:T}\end{equation}

and can be approximated by the score-function estimator described in \citep{poyiadjisdoucetsingh2011particlescore}. Since complete ancestral trajectories are available from the forward filter, we use

\begin{equation}
    \widehat{S}_{N,T}(\Theta) = \sum_{i=1}^N w_T^{(i)} \nabla_\Theta \log p(x_{0:T}^{(i)}, y_{1:T} \mid \Theta)
\end{equation}

where $w_T^{(i)}$ are the (normalized) particle weights at the final time $T$ and $\widetilde{x}_{0:T}^{(i)}$ is the complete ancestral lineage of terminal particle $i$. This estimator is consistent under standard assumptions but generally biased for finite $N$. 

% For the Rao--Blackwellized filter, whose particles are Gaussian components rather than complete paths, the complete-data term is replaced by its conditional expectation given the Rao--Blackwellized lineage,

% \begin{equation}
%     \widehat{S}_{N,T}^{\mathrm{RB}}(\Theta) = \sum_{i=1}^N w_T^{(i)} \mathbb{E}\left[ \nabla_\Theta \log p(X_{0:T}, y_{1:T} \mid \Theta) \,\middle|\, \text{RB lineage } i, y_{1:T} \right]
% \end{equation}

% or, equivalently, full trajectories are first generated using the RB backward sampler and then inserted into the preceding path-space approximation. 

Furthermore, we can estimate the log-likelihood gradient using the stop-gradient trick described in \citep{scibior2021differentiableparticlefiltering} without modifying our filter. With resampling at every step and carried numerical weights $1/N$, the normalizing-constant estimator is

\begin{equation}
    \widehat Z_0 = 1, \qquad \widehat Z_t = \widehat Z_{t-1} \left[ \frac{1}{N} \sum_{i=1}^N G_t(y_t \mid x_t^{\mathcal{O}_t,(i)}) \right], \nonumber
\end{equation}

so that
\begin{equation}
    \log \widehat Z_T = \sum_{t=1}^T \log \left[ \frac{1}{N} \sum_{i=1}^N G_t(y_t \mid x_t^{\mathcal{O}_t,(i)}) \right]. \nonumber
\end{equation}

% For stop-gradient resampling, define $\operatorname{sg}(z)$ as numerically equal to $z$ with zero derivative. If $b_i$ is the selected ancestor, the corrected post-resampling surrogate log-weight is

% \begin{equation}
%     \ell_{t,\mathrm{res}}^{(i)} = -\log N + \log w_{t-1}^{(b_i)} - \operatorname{sg}\left( \log w_{t-1}^{(b_i)} \right), \nonumber
% \end{equation}

% whose forward value is $-\log N$ but whose derivative retains the resampling score, and the next surrogate unnormalized log-weight is
% \begin{equation}
%     \widetilde{\ell}_t^{(i)} = \ell_{t,\mathrm{res}}^{(i)} + \log G_t(y_t \mid x_t^{\mathcal{O}_t,(i)}). \nonumber
% \end{equation}

% The corrected claim is that automatic differentiation of the surrogate yields a consistent estimator of the score,
% \begin{equation}
%     \widehat{S}_{N,T}^{\mathrm{SGR}}(\Theta) := \operatorname{AD}_\Theta\left[ \log \widehat Z_T^{\mathrm{SGR}} \right] \xrightarrow[N \to \infty]{\mathbb{P}} \nabla_\Theta \log p(y_{1:T} \mid \Theta), \nonumber
% \end{equation}

% under the regularity assumptions of the particle score estimator and with the proposal-gradient treatment required by the stop-gradient-resampling construction. This is a consistent estimator of the score, not an exact equality or a generally finite-$N$ unbiased estimator.

\section{Performance comparison}

Evaluate the particle diversity during the run.

Pre-resampling covariance

\begin{equation}
    m_{pre} = \sum_{i=1}^N w_t^{(i)} m_t^{(i)}
\end{equation}

\begin{equation}
    C_{pre} = \sum_{i=1}^N w_t^{(i)} \left( m_t^{(i)} - m_{pre} \right) \left( m_t^{(i)} - m_{pre} \right)^T
\end{equation}

Post-resampling covariance

\begin{equation}
    m_{post} = \sum_{i=1}^N \frac{1}{N} m_t^{a(i)}
\end{equation}

\begin{equation}
    C_{post} = \sum_{i=1}^N \frac{1}{N} \left( m_t^{a(i)} - m_{post} \right) \left( m_t^{a(i)} - m_{post} \right)^T
\end{equation}
